\section{Polynomial multiplication near contact}\label{sec:profiles}

The tripod principle compares the price $D$ of a common label at the
three terminals with the lower bound $D_*(a)=\sum_i a_i-(\sum_i a_i^{-1})^{-1}$
for their arm lengths $a_i>0$. We call equality \emph{contact} and write
$\Delta=D-D_*(a)$ for the defect. To propagate contact to other triples
of arm lengths, we need to know how a price changes when just one input
width varies. Polynomial multiplication supplies a tensor for which we
can measure this change. Its price as one input grows has a limiting slope
and intercept after exponentiation. Proposition~\ref{thm:profile-rigidity}
determines the intercepts to within $O(\Delta)$. Section~\ref{sec:energy}
then compares different placements, and Section~\ref{sec:contact} uses the
comparison to propagate equality.

Throughout these three sections the field is algebraically closed and
prices use natural logarithms with the distances held fixed. A label
with $q$ values shared at the terminals costs $D\log q$. On an embedded
tripod, including its origin gives the same price. A duplicate at either
of two terminals is free, and convexity bounds its cost at every point
on the geodesic between them.

\subsection{The continuum profile}
Place a finite nonempty set of inputs and one output in a spectral space.
Write $H$ for price and $D>0$ for the common-label rate of the placement.
For integer lengths $L_i\ge1$, let $M(L)$ multiply univariate polynomials
of degrees less than $L_i$, retaining all output coefficients. The following
normalization also fixes the constant in the energy formula below.

\begin{lemma}\label{lem:continuum-profile}
There is a continuous coordinatewise nondecreasing function $P$ on the positive
orthant such that, if $L_i/N\to t_i>0$, then
\[
 H(M(L))=D\log N+P(t)+o(1),\qquad P(ct)=P(t)+D\log c.
\]
For polynomials in $m$ variables with total input degrees at most
$\lfloor Nt_i\rfloor$, the corresponding price is
$mD\log N+mP(t)-D\log(m!)+o(1)$ for each fixed $m$.
\end{lemma}

\begin{proof}
Evaluation and interpolation bound $H(M(L))$ by $D\log(\sum_iL_i)$.
In the other direction, $M(kL)$ degenerates to $k$ direct copies of
$M(L)$. Indeed, divide $[0,1]$ into $k$ equal intervals and choose a
continuous convex piecewise affine function with a strict change of slope
at each division. An input monomial of degree $d_i$ has normalized position
$(d_i+1/2)/(kL_i)$. Give it weight $kL_i$ times the function there. Give the
output the opposite weight at the weighted mean of those positions, with
coefficient $k\sum_i L_i$. The output weight depends only on the sum of
degrees. Jensen's inequality retains precisely the terms whose inputs lie
in one interval. Different intervals use disjoint input and output
coordinates, and each surviving block is a translated copy of $M(L)$.
The weight equalities define a rational linear subspace. Rational
approximation within it preserves the finitely many strict inequalities,
and scaling gives integral weights. Rescaling the coordinates and using
closedness of price sublevels therefore gives the asserted degeneration.

For $L_i=\lfloor Nt_i\rfloor$ the normalized cost at $N$ therefore
dominates that at any fixed sufficiently large $M$, up to
$D\log(M\lfloor N/M\rfloor/N)$. Taking first the lower limit in $N$ and
then the upper limit in $M$ proves convergence. Changing the unit of $N$
gives log homogeneity, first with floors by squeezing the scales and then
for real $N$. Coordinate monotonicity and squeezing between
$(1\pm\epsilon)t$ give continuity and allow $o(N)$ errors in the lengths.

For the assertion about total degree, divide the exponent vector at
input $i$ by $Nt_i$ and at the output by $N\sum_i t_i$. All inputs now
range over lattice points in the unit simplex. To cut this domain by a
hyperplane, first discard input points on the plane. Apply the same
convex weighting as above, using the positive part of the defining
affine form and coefficients $Nt_i$. Jensen's inequality retains terms
whose inputs are on the same side. Their outputs are on that side as
well, so the two blocks form a direct sum. One may repeat the cut within
each block.

On a translated axis-parallel cube of side $s$, multiplication factors
over the $m$ coordinates. The univariate limit shows that its
exponentiated price, divided by $N^m\exp(mP(t)/D)$, tends to its volume
$s^m$. Translation just relabels exponents and rounding a cube side
changes widths by $O(1)$. Packing the simplex with interior grid cubes
thus bounds its normalized price from below by $1/m!$, since
$\exp(H/D)$ adds on direct sums. For the upper bound, start with an
enclosing cube and cut off a slightly enlarged simplex, so all input
points of the desired tensor survive in one cell. Pack the complementary
cells with interior grid cubes. The enclosing tensor dominates the
direct sum of the desired tensor and all packed cubes. Refining the
grid and shrinking the enlargement accounts for the whole complementary
volume and gives the matching upper bound $1/m!$. Taking logarithms
proves the formula.
\end{proof}

\subsection{A label visible at two sites}

We also need a lower bound for a sum whose block label is visible at
only two of three sites. Suppose $T=\sum_b T_b$ with $T_b\ne0$ and
$T_b\in X\otimes Y_b\otimes Z_b$, where the $Y_b$ are direct subspaces
at $Y$, and the $Z_b$ are direct subspaces at $Z$. For every probability
law $q$ on the block labels, writing
$\operatorname{Ent}(q)=-\sum_b q_b\log q_b$, we have
\begin{equation}
 H(T)\ge\sum_b q_b H(T_b)+d(Y,Z)\operatorname{Ent}(q).
 \label{eq:pair-entropy}
\end{equation}
To see the entropy term, take a rational $q$ and a power $n$ for which
each $nq_b$ is integral. At $Y$ and $Z$, retain the
$M=n!/\prod_b(nq_b)!$ words of type $q$. Each word has price
$n\sum_bq_bH(T_b)$. The shared-leg separation of~\cite[Proposition~3.1]{nine},
in the interpolation form of~\cite[Proposition~2.1]{lu}, degenerates
$2M-1$ direct copies of the retained tensor to the full direct sum of
the $M$ word blocks, each tensored with an $M$-dimensional wire on
$Y,Z$. Interpolation uses $2M-1$ distinct field elements, available
in every algebraically closed field. Comparing the two prices gives
\[
 nH(T)+D\log(2M-1)\ \ge\
 n\sum_bq_bH(T_b)+(D+d(Y,Z))\log M ,
\]
where $D$ is the common-label rate of the three sites.
Now $\log((2M-1)/M)\le\log2$ and $n^{-1}\log M\to\operatorname{Ent}(q)$.
Division by $n$ and continuity in $q$ give~\eqref{eq:pair-entropy}.

\subsection{Slopes and intercepts on a triangle}

Use $i,j,k$ for the three sites in any order and suppose their arms $a_i$
are positive. Put $d_i=d(j,k)=a_j+a_k$, $D=\delta(\{i,j,k\})$, and
write $P_k(x,y)$ for the profile with output $k$ and input widths $x$ at
$i$ and $y$ at $j$. We identify each argument by its site.
For a fixed first input define
\[
 F_{ik}(t)=\exp\bigl(P_k(t_i=1,t_j=t)/d_i\bigr).
\]
The determinant filtration and the shifted three-block degeneration for
polynomial multiplication \cite[proofs of Lemmas~4.1 and~4.2]{nine} give
\[
 F_{ik}\text{ is concave},\qquad
 F_{ik}(3t+1)\ge 2F_{ik}(t)+F_{ij}(t).
\]
In the first degeneration a two-dimensional wire on $j,k$ supplies two
blocks, with the length at $j$ increased and decreased by one. The blocks
are direct on that pair. Optimizing the pair-label inequality over the
block probabilities adds their exponentiated prices at scale $d_i$.
This gives midpoint concavity on the lattice of lengths, hence concavity
in the limit. In the second degeneration the longer length is
$3L_j+L_i-1$, and one of the three blocks exchanges $j$ and $k$. The same
pair-label inequality and Lemma~\ref{lem:continuum-profile} give the second
displayed assertion. Both degenerations work in every characteristic.
For the first, injectivity of multiplication by the determinant and
a dimension count give the degree-one Clebsch--Gordan determinant
sequence (see~\cite{kowalski}). Lifting monomials to its filtration
suffices. The second uses a monomial grading, and reversal of the order
on a leg identifies the middle block.

Positive increasing concave functions have finite limiting slopes at
infinity. Dividing the three-block inequalities by $t$ shows that the
slopes of $F_{ik}$ and $F_{ij}$ are equal. For $S=F_{ik}+F_{ij}$ we have
$S(3t+1)\ge3S(t)$. Iterating along
$t_n=3^n(t+1/2)-1/2$ gives
$S(t)/(t+1/2)\le \lim_{n\to\infty}S(t_n)/(t_n+1/2)$, twice the common slope. In particular the common slope is positive. Write it as
$\exp(c_i/d_i)$.
For either function the excess over the line through zero of that slope
is nonnegative and increasing. The bound on $S$ therefore gives numbers
$\beta_{ik},\beta_{ij}\ge0$, with $\beta_{ik}+\beta_{ij}\le1$, such that
\[
 t\le e^{-c_i/d_i}F_{ik}(t)\le t+\beta_{ik}.
\]
Here $\beta_{ik}$ is the limiting excess after division by the slope.
In particular $P_k$ is $d_i$-Lipschitz as a function of $\log y$ with $x$
fixed, since every one-sided derivative of $F_{ik}$ lies between zero and
$F_{ik}(t)/t$.

Set $h=(\sum_i1/a_i)^{-1}$, $b_i=a_i-h$, $u_i=b_i/a_i$,
$\beta_{ik}^*=a_k/d_i$, and $\Delta=D-D_*(a)$, where
$D_*(a)=\sum_i a_i-h$. Since $P_k(x,y)=D\log x+P_k(1,y/x)$
and $D-d_i=b_i+\Delta$, the intercept bound gives an upper envelope
valid at every pair of widths,
\begin{equation}\label{eq:profile-upper}
 P_k(x,y)\le U_{ik}(x,y):=c_i+(b_i+\Delta)\log x+d_i\log(y+\beta_{ik}x).
\end{equation}

\begin{proposition}[Profile near contact]\label{thm:profile-rigidity}
We have $\Delta\ge0$. If $\min_i u_i\ge\eta>0$, there is a number $C$
common to all three sites such that
\begin{equation}\label{eq:src2}
 d_i(\beta_{ik}-\beta_{ik}^*)=O_\eta(\Delta),\qquad
 c_i+b_i\log u_i=C+O_\eta(\Delta).
\end{equation}
In the sector $y/x\ge u_j/u_i$,
\begin{equation}\label{eq:src3}
 0\le U_{ik}(x,y)-P_k(x,y)\le O_\eta(\Delta).
\end{equation}
\end{proposition}

\begin{proof}
Let $E_{ik}=U_{ik}-P_k$ and $z_{ik}=d_i(\beta_{ik}-\beta^*_{ik})$.
Differentiating with respect to $y$ and bounding the derivative of
$F_{ik}$ below by its limiting slope yields
\[
 \partial_y P_k(x,y)\ \ge \frac{d_i}{y+\beta_{ik}x}\exp(E_{ik}(x,y)/d_i).
\]
There is a second bound from interchanging $i,j$. At differentiability
points, $x\partial_xP_k+y\partial_yP_k=D$. The profiles are locally
Lipschitz, so continuity gives the sum of the two derivative bounds at
$(x,y)=(u_i,u_j)$ even if a derivative fails to exist there.
At this point $u_j+u_i\beta^*_{ik}=1$ and $u_jd_i+u_id_j=D_*$.
Use $(1+r)^{-1}=1-r+r^2/(1+r)$ with
$r=u_i z_{ik}/d_i$, and $e^s\ge1+s$ for $s\ge0$.
Since $\eta\le u_j+u_i\beta_{ik}\le2$, this gives, for some $c_\eta>0$,
\[
 \frac{\Delta}{u_iu_j}\ge -(z_{ik}+z_{jk})+
 c_\eta\left(E_{ik}+E_{jk}
           +\frac{z_{ik}^2}{d_i}+\frac{z_{jk}^2}{d_j}\right),
\]
where all the gaps are evaluated at the specified widths.
The row sums $z_{ik}+z_{ij}$ are nonpositive since
$\beta^*_{ik}+\beta^*_{ij}=1$. Summing over output sites $k$ gives
$\Delta\ge0$ and bounds each nonnegative remainder by $O_\eta(\Delta)$.
Each row sum thus lies between $-O_\eta(\Delta)$ and zero. Each
column sum is at least $-O_\eta(\Delta)$ by the displayed inequality.
Since the columns have the same total as the rows, each column sum is
also at most $O_\eta(\Delta)$.

Comparing the two envelopes for $P_k(u_i,u_j)$ now gives
\[
 c_i+b_i\log u_i+u_i z_{ik}
 =c_j+b_j\log u_j+u_j z_{jk}+O_\eta(\Delta).
\]
Here the logarithmic Taylor remainders are controlled by
$z_{ik}^2/d_i$ and $z_{jk}^2/d_j$. The row and column estimates
show that, up to $O_\eta(\Delta)$, the entries
$(z_{12},z_{13},z_{23},z_{21},z_{31},z_{32})$ alternate between
$z_{12}$ and $-z_{12}$. Sum the three comparisons of envelopes for
$(i,j)=(1,2),(2,3),(3,1)$. The terms $c_i+b_i\log u_i$ cancel, leaving
$2(u_1+u_2+u_3)z_{12}=O_\eta(\Delta)$.
This proves \eqref{eq:src2}. The gap $E_{ik}$ depends only on $y/x$. The
difference between $t+\beta_{ik}$ and $e^{-c_i/d_i}F_{ik}(t)$ decreases
with $t$, and the latter function increases. Their log ratio hence
decreases, which proves \eqref{eq:src3}.
\end{proof}

For example, at contact on the unit tripod the predicted intercept
is $\beta^*_{ik}=1/2$ and $u_i=u_j=2/3$. Thus
$e^{-c_i/d_i}F_{ik}(t)=t+1/2$ for $t\ge1$. More generally,
\eqref{eq:src2} remains uniform as a single arm tends to infinity
while the other two stay in a compact subset of $(0,\infty)$.
If an arm has length zero, $D\ge D_*(a)$ follows from the diameter bound.

\section{Energy and radial contraction}\label{sec:energy}

Fix an output and a nonempty finite set $A$ of inputs with common-label
rate $D>0$ on all the terminals. For costs $\ell_i>0$ per unit of input
width, define the energy
\begin{equation}\label{eq:energy-widths}
 Q_A(\ell)=D+\sup_{t_i>0}\bigl(P(t)-\ell\cdot t\bigr).
\end{equation}
It is finite by the evaluation bound, convex in $\ell$, and satisfies
$Q_A(c\ell)=Q_A(\ell)-D\log c$ by log homogeneity of $P$.
A common surcharge $s$ penalizes a candidate width vector by
$s\sum_i t_i$. Hence the slope of the optimized energy along a common
cost increase measures total width. We will use the finite difference
$Q_A(\ell)-Q_A(\ell+s\mathbf1)$.
Lemma~\ref{lem:energy-compare} shows that this drop can only decrease
when inputs move towards the output and their respective costs decrease
in proportion. We prove the comparison first for homogeneous degrees,
where we can move an input between two polynomial multiplications,
and then recover $Q_A$ by rescaling those degrees.

To compare different placements, first charge a cost $f_i>0$ per
degree of a homogeneous polynomial.
Write $M_m^{\mathrm{hom}}(r)$ for multiplication of homogeneous
polynomials in $m$ variables, with input degrees $r=(r_i)_{i\in A}$.
Define
\[
 B_A(f)=\sup_{m\ge1,\ r_i\in\mathbb Z_{\ge0}}
 \frac{H(M_m^{\mathrm{hom}}(r))-\sum_i f_i r_i}{m}.
\]
This is finite. Evaluation and interpolation at a basis of evaluation
functionals on the output give rank at most its dimension. With
$R=\sum_i r_i$ and $q=\exp(-\min_i f_i/D)$, the exponential of the
cost-subtracted price at scale $D$ is bounded by
$\binom{m-1+R}{R}q^R\le(1-q)^{-m}$. The supremum is unchanged if $m$
is required to be a multiple of a prescribed integer. Indeed, using
disjoint groups of variables recovers any power of the original tensor
as a restriction of multiplication with the combined degrees.

\begin{lemma}[Adding an input]\label{lem:increasing-returns}
If $\varnothing\ne A'\subseteq A$ and $s>0$, then
\begin{equation}\label{eq:src5}
 B_A(f)-B_A(f+s\mathbf1)\ge
 B_{A'}(f)-B_{A'}(f+s\mathbf1),
\end{equation}
with costs restricted to the indicated input set. If all sites in a
restricted placement coincide, its energy is zero.
\end{lemma}

\begin{proof}
It suffices to add a single input $e$. Rewrite the claim with
$B_{A'}(f)+B_A(f+s\mathbf1)$ on the right, and choose two multiplications
approaching those suprema. They can use the same number $m$ of variables
by taking powers in separate groups. Name their variables respectively
$u=(u_1,\ldots,u_m)$ and $v=(v_1,\ldots,v_m)$. At each common input the
tensor product of the two multiplications accepts a bihomogeneous
polynomial in $(u,v)$ of prescribed bidegree. At $e$ it accepts a
polynomial $h(v)$ of degree $q$. Write $(U,V)$ for the total bidegree,
including $q$.

We realize this tensor product by a sum of multiplications in which $e$
belongs instead to the $u$ group. Rename $h(v)$ to $h(u)$ and substitute
$u-v$ for $u$ at the common inputs. Multiply, substitute $u+v$ for $u$
on the output, and project onto bidegree $(U,V)$. The substitutions on
the common inputs cancel, and the only term of $h(u+v)$ that survives
the projection is $h(v)$. This identity uses no division. Expand the
intermediate inputs into bidegrees before multiplying. If a total of $k$
degrees from the common inputs move from $u$ to $v$, the intermediate
$v$-degree is $V-q+k$. The final substitution can only increase this
degree, so a surviving sector has $k\le q$. The total degree at each
input is preserved. Thus the base cost of every sector is unchanged and
its surcharge on $v$ is at most $sV$. Each sector factors through a
multiplication on $u$ with all inputs and one on $v$ with inputs $A'$,
with local pre- and postprocessing.

Apply this construction to powers in $nm$ variables per group with all
degrees multiplied by $n$. There are polynomially many bidegree sectors
as functions of $n$, since $A$ is fixed. The price bound for an ordinary
sum costs only $O(\log n)$ beyond the largest sector price. After division
by $nm$ and passage to the limit, the right side is bounded by
$B_A(f)+B_{A'}(f+s\mathbf1)$, as claimed.
\end{proof}

\subsection{From degrees to widths}

Charging $\ell_i/N$ per degree makes degrees of order $N$ correspond
to widths of order one. The factor $1/m!$ for total-degree polynomials
in Lemma~\ref{lem:continuum-profile} accounts for the additive $D$ in
\eqref{eq:energy-widths}.

\begin{lemma}[Continuum limit of the degree energy]\label{lem:continuum-energy}
At a placement with common-label rate $D>0$, one has
\begin{equation}\label{eq:src6}
 B_A(\ell/N)=D\log N+Q_A(\ell)+o(1)
 \qquad(N\longrightarrow\infty).
\end{equation}
\end{lemma}

\begin{proof}
Write $S_m(r)$ for polynomial multiplication with total degree
\emph{at most} $r_i$ on each input. Its price minus degree costs, divided
by $m$, is a lower bound on $B_A$. Indeed, $n$ copies on disjoint variable
groups restrict from homogeneous multiplication with $nm+1$ variables
and degrees $nr_i$. Use the last variable to homogenize and set it to one
on the output. Hence
$B_A(f)\ge n(H(S_m(r))-f\cdot r)/(nm+1)$, and we let $n\to\infty$.
For $r_i=\lfloor Nm t_i\rfloor$, Lemma~\ref{lem:continuum-profile} now gives
\[
 \liminf_{N\to\infty}(B_A(\ell/N)-D\log N)
 \ge P(t)-\ell\cdot t+D\bigl(\log m-m^{-1}\log(m!)\bigr).
\]
The last summand tends to $D$. Taking the supremum over $t$ gives the
required lower bound.

For the upper bound take integer $N$. In the supremum defining $B_A$
we may use $Km$ variables grouped into $K$ sets of size $m$. At each
input, record the total degree in each group by its bin
$[(b_i-1)N,b_iN)$, $b_i\ge1$. Within a fixed bin pattern the
multiplication restricts from the product of $S_m(Nb)$ over the groups,
and the cost in each group is at least $\sum_i\ell_i(b_i-1)$. The price
of the sum over patterns is bounded by the logarithm of the sum of their
exponentiated prices at scale $D$. Summing over all bins independently
in each group enlarges this last sum to a $K$th power. Dividing its
logarithm by $Km$ and taking the supremum gives
\[
 B_A(\ell/N)\le \frac Dm\log\sum_{b\in\mathbb Z_{>0}^A}
 \exp\left(\frac{H(S_m(Nb))-\sum_i\ell_i(b_i-1)}{D}\right).
\]
For fixed $m$, after removing $N^m$ from the sum, the terms are dominated
by a constant times $(m+\sum b_i)^m\exp(-\ell\cdot b/D)$, by the
output-dimension bound. Lemma~\ref{lem:continuum-profile} and dominated
convergence bound the upper limit of $B_A(\ell/N)-D\log N$ by
\[
 D\left(\log m-\frac{\log(m!)}m\right)
 +\frac{\sum_i\ell_i}m+
 \frac Dm\log\sum_b\exp\left(\frac{mP(b/m)-\ell\cdot b}{D}\right).
\]
For $0<\epsilon<1$, use the supremum of
$P(t)-(1-\epsilon)\ell\cdot t$ to bound the last term, with a remainder
at most $(D/m)\log\sum_b\exp(-\epsilon\ell\cdot b/D)$.
Let $m\to\infty$, then $\epsilon\downarrow0$, using the scaling relation
for $Q_A$. Monotonicity in $N$ gives the same limit at real scales.
\end{proof}

\begin{lemma}[Radial contraction]\label{lem:energy-compare}
Move each input $i$ towards the output along its geodesic, to a fraction
$0<\lambda_i\le1$ of its original distance. Write $A'$ for the resulting
placement and $\ell'_i=\lambda_i\ell_i$. For $s>0$ one has
\begin{equation}\label{eq:src7}
 Q_A(\ell)-Q_A(\ell+s\mathbf1)\ge
 Q_{A'}(\ell')-Q_{A'}(\ell'+s\mathbf1)
\end{equation}
when both placements have positive common-label rate.
\end{lemma}

\begin{proof}
Duplicating an input at its own site and cost leaves $B$ unchanged, by
local multiplication. An additional input at the output at zero cost
also leaves $B$ unchanged, by local multiplication on the output.
The reverse inequalities in both statements use degree zero at the
added input. For each fixed multidegree the price is convex as the
added input moves on a geodesic, and the subtracted cost is affine.
Thus an added input at radial fraction $\lambda_i$ and cost at least
$\lambda_i f_i$ leaves $B$ unchanged, first for one addition and then
for all of them. This applies both at $f=\ell/N$ and after increasing
every cost by $s/N$, since $\lambda_i\ell_i+s\ge
\lambda_i(\ell_i+s)$. Insert the moved inputs at their specified costs
and use Lemma~\ref{lem:increasing-returns} to delete the original
inputs. The energy drop for the resulting placement is bounded by the
drop for the original one. Let $N\to\infty$ in \eqref{eq:src6}. The
leading $D\log N$ cancels separately at each placement, which proves
\eqref{eq:src7}.
\end{proof}

\section{Propagation of equality on a tripod}\label{sec:contact}

Fix three infinite rays forming a metric tripod with origin $O$.
Prices on these rays are homogeneous under simultaneous radial
dilation. Proposition~\ref{thm:profile-rigidity} gives
$D(a)\ge D_*(a)$ at the three tips of lengths $a_1,a_2,a_3$.
We now prove the rigidity assertion of Proposition~\ref{prop:tripod}.
Equality at a single triple of lengths forces equality at every
triple. The same conclusion follows from equality approached along
one arm as its length tends to infinity with the other two lengths
fixed and positive.

Use the reciprocal lengths
\[
 v_i=1/a_i,\qquad h(a)=\frac{1}{v_1+v_2+v_3},\qquad
 \Delta(v)=D(a)-D_*(a).
\]
An infinitely long arm now corresponds to a coordinate face.
The proof first compares integrals of the defect along the three
coordinate rays out of $v$. At a zero of $\Delta$ these integrals
agree. Radial homogeneity then forces their common value to vanish.

\subsection{Comparing the three directional integrals}

The profile lower bound and the spanning tree through $O$ give
$D_*(a)\le D(a)\le\sum_i a_i$. Equivalently,
\begin{equation}\label{eq:defectbound}
 0\le \Delta(v)\le\frac{1}{v_1+v_2+v_3},\qquad
 \Delta(rv)=r^{-1}\Delta(v)\quad(r>0).
\end{equation}
On the positive octant $\Delta$ is locally Lipschitz and locally
semiconvex. The latter means that adding $C|v|^2$ makes it convex
on a given small ball for a suitable $C$. Indeed, $D$ is convex in
the arm lengths by joint geodesic convexity. Transport of a single
tip through distance $t$ changes $D$ by at most $t$. The change of
variables $a_i=1/v_i$ has bounded first and second derivatives on
interior compact sets, and $D_*$ is smooth there.

As $v_1\downarrow0$, the function $D(a)-\sum_i a_i$ is
nonincreasing by the same transport bound, and is bounded below by
$-h(a)$. Since $h(a)$ also has a limit there, $\Delta$ extends to the open face $v_1=0$.
The Lipschitz and semiconvexity constants in the other two variables
are uniform when those variables range in compact subsets of
$(0,\infty)$, including as $v_1\downarrow0$. Thus the extension is
continuous and has the same tangential regularity. The other faces
are treated by relabeling.

For each interior $v$ and coordinate vector $e_i$, set
\[
 J_i(v)=\int_0^\infty
 \left[\Delta(v+t e_i)-\mathbf 1_{\{t\le1\}}\Delta(v)\right]
       \frac{dt}{t}.
\]
The subtraction makes the integral convergent at zero by local
Lipschitz continuity. Equation~\eqref{eq:defectbound} gives
convergence at infinity. We will prove the estimate
\begin{equation}\label{eq:src10}
                   |J_i(v)-J_j(v)|\le C\Delta(v).
\end{equation}
Here $C$ is locally uniform up to each open face. On a face the same
estimate holds for its two tangential directions.

\paragraph{An energy derivative.}
Keep $v$ fixed and use $a_i=1/v_i$ as the locations of three input
tips, with the output at $O$. Let $Q$ be their energy from
Equation~\eqref{eq:energy-widths}. For a variable
$s\in(0,\infty)^3$, put
\[
 G_v(s)=Q\big((a_i s_i)_i\big)+\sum_i a_i\log s_i
                   -h(a)\log\Big(\sum_i s_i\Big).
\]
Moving $s$ in direction $v$ raises every cost $a_i s_i$ at the same
unit rate. On the same tripod place tip $i$ at distance $r/s_i$
with cost $r$ on every input. The ratio of cost to distance on arm
$i$ is still $s_i$. For small $r$ the comparison placement lies
inside the fixed one, and for large $r$ it lies outside. By
Lemma~\ref{lem:energy-compare}, the energy drops at these inner and
outer placements bound the drop at the fixed placement from below
and above, respectively.

At the comparison placement the common-label rate is
$rD((1/s_i)_i)$. Adding $O$ to the three label sites costs nothing.
Increasing all costs from $r$ to $r+t$ therefore lowers the energy
by $rD((1/s_i)_i)\log(1+t/r)$, by the scaling law for $Q$.
After division by $t$, both bounds tend to $D((1/s_i)_i)$ as
$t\downarrow0$. Differentiating the two explicit logarithmic terms
of $G_v$ adds $\sum_i1/s_i-1/\sum_i s_i=D_*((1/s_i)_i)$. Hence
\begin{equation}\label{eq:src8}
       \left.\frac{d}{dt}\right|_{t=0+}G_v(s+tv)=-\Delta(s).
\end{equation}
The finite convex function $Q$ is locally Lipschitz, so this identity
may be integrated along any ray in direction $v$.

\paragraph{Starting the rays near their axes.}
Let $s_i=1$, $s_j=s_k=\varepsilon$. Recall $b_i=a_i-h(a)$ and the
intercept $c_i$ of Section~\ref{sec:profiles}. Write
$B_i=b_i+\Delta(v)$ and $d_i=a_j+a_k$. The limit of $G_v$ is
\begin{equation}\label{eq:src9}
           \lim_{\varepsilon\downarrow0}G_v(s)
                =c_i+B_i\log\frac{B_i}{a_i}.
\end{equation}
To prove the upper bound, let the three input widths be $x,t_j,t_k$
and denote their profile with output at $O$ by $P(x,t_j,t_k)$.
The point $O$ divides the segment from $j$ to $k$ in the ratio
$a_j:a_k$. Convexity in the output position bounds $P$ by the
weighted average of the profiles with output at $k$ and at $j$,
with respective weights $a_j/d_i$ and $a_k/d_i$. At either
endpoint, the input at the output site can be multiplied in
locally. Equation~\eqref{eq:profile-upper} and
$\beta_{ik},\beta_{ij}\le1$ give
\[
 P(x,t_j,t_k)\le c_i+B_i\log x
                +a_j\log(t_j+x)+a_k\log(t_k+x).
\]
In the energy formula use costs $(a_i,\varepsilon a_j,\varepsilon a_k)$.
Changing variables to $x,t_j+x,t_k+x$ makes the effective cost of
$x$ equal to $a_i-\varepsilon d_i>0$ for small $\varepsilon$.
Maximizing the three terms separately and using $B_i+d_i=D(a)$
gives
\[
 Q\le c_i+B_i\log\frac{B_i}{a_i-\varepsilon d_i}
                  +d_i\log(1/\varepsilon).
\]

For a matching lower bound take $x=B_i/a_i$, $t_j=1/\varepsilon$,
$t_k=x+t_j$. At $O$ extract a single coefficient of the product.
Choosing complementary degrees and reversing the basis at $k$ makes
this coefficient the pairing of the third input with the product
of the first two. An endpoint rounding by one leaves the continuum
profile unchanged. Hence $P(x,t_j,t_k)\ge P_k(x,t_j)$, and the
asymptotic slope from Section~\ref{sec:profiles} gives
\[
       P_k(x,t_j)=c_i+B_i\log x+d_i\log(1/\varepsilon)+o(1).
\]
The total width cost of this choice tends to $B_i+d_i=D(a)$, which
cancels the additive constant in the energy formula. Subtracting
$d_i\log(1/\varepsilon)$ through the logarithmic terms of $G_v$
gives Equation~\eqref{eq:src9}.

\paragraph{Comparison by integration.}
Start two rays near $e_i$ and $e_j$, at points of the form used in
Equation~\eqref{eq:src9}. Integrate Equation~\eqref{eq:src8} in
direction $v$ on both rays and subtract. Their terminal $G_v$ values
have difference tending to zero at infinity, since
$G_v(rs)=G_v(s)-\Delta(v)\log r$ and $G_v$ is continuous at $v$.
Now send the starting points to $e_i$ and $e_j$. The two integrands
are uniformly bounded near zero by Equation~\eqref{eq:defectbound}.
At infinity their difference is $O(t^{-2})$ by homogeneity and the
Lipschitz bound near $v$. The difference of the two axis limits in
Equation~\eqref{eq:src9} therefore equals
\[
 \int_0^\infty\big[\Delta(e_i+tv)-\Delta(e_j+tv)\big]\,dt
    =J_i(v)-J_j(v).
\]
For the last identity use $\Delta(e_i+tv)=t^{-1}\Delta(v+t^{-1}e_i)$
and substitute $u=1/t$.

By Equation~\eqref{eq:src2}, the quantities $c_i+b_i\log(b_i/a_i)$
agree up to $O(\Delta(v))$. Replacing $b_i$ by $B_i$ changes each
quantity by $O(\Delta(v))$ as well. Indeed, $b_i/a_i$ is bounded
away from zero when the other two reciprocal lengths stay in compact
subsets of $(0,\infty)$ and $v_i$ approaches zero; more generally,
all three ratios stay bounded away from zero in a neighborhood of any
open face point. Also $B_i/a_i\le1$ by Equation~\eqref{eq:defectbound}.
The derivative of $b\mapsto b\log(b/a_i)$ is uniformly bounded
throughout the interval in question. This proves Equation~\eqref{eq:src10}
with the asserted uniformity. For the two tangential directions
at a face, dominated convergence uses the tangential Lipschitz bound
at zero and Equation~\eqref{eq:defectbound} at infinity.

\subsection{The strong minimum principle}

\begin{proof}[Completion of the proof of Proposition~\ref{prop:tripod}]
Suppose $\Delta(v)=0$. Work either in the interior or on an open face,
using just its positive coordinates, indexed $1,\ldots,k$ after
relabeling. All $J_i(v)$ equal a common $J\ge0$ by
Equation~\eqref{eq:src10}. A simultaneous dilation gives the exact
change $J_i((1+\varepsilon)v)=J/(1+\varepsilon)$. To compare this
decrease with the variation on the individual rays, perform the
dilation one coordinate at a time. Its intermediate points are
\[
          v^i=v+\varepsilon\sum_{j\le i}v_j e_j
               \qquad(0\le i\le k).
\]
At each intermediate point the cost of switching from $J_i$ to
$J_{i+1}$ is bounded by $C\Delta(v^i)$. Thus
\[
 S_\varepsilon=\sum_{i=1}^k\big(J_i(v^i)-J_i(v^{i-1})\big)
   =\frac{J}{1+\varepsilon}-J+
                  O\!\left(\max_{0<i<k}\Delta(v^i)\right).
\]
Write $K=\limsup_{\varepsilon\downarrow0}
\max_{0<i<k}\Delta(v^i)/\varepsilon$, finite by local Lipschitz
continuity, and set $g_i(t)=\Delta(v+te_i)$. We claim that
\begin{equation}\label{eq:src11}
 -J+CK\ \ge\ \liminf_{\varepsilon\downarrow0}
              \frac{S_\varepsilon}{\varepsilon}
        \ \ge\ \sum_{i=1}^k v_i\int_0^\infty g_i'(t)\,\frac{dt}{t}.
\end{equation}

To justify the second inequality, write all $k$ integrals with the
same variable $t$. For almost every $t>0$, the increment of $\Delta$
on $[v^{i-1}+te_i,v^i+te_i]$, divided by $\varepsilon$, tends to
$v_i g_i'(t)$. Indeed, after adding a local quadratic, the slope
on this segment lies between the slopes of adjoining segments in
direction $e_i$ of any fixed small length. Send $\varepsilon$ to
zero, then shrink the adjoining segments at each point where $g_i$
is differentiable. The unshifted increments that are subtracted for
$t\le1$ sum to $\Delta((1+\varepsilon)v)-\Delta(v)=0$.

These subtractions also control the singularity at zero. Pair each
shifted increment with its own unshifted increment. Adding a local
quadratic makes the shifted increment at least the unshifted one,
because both segments have the same length and direction. Removing
the quadratic gives a lower bound $-C\varepsilon t$ for the
difference. For large $t$, moving the starting point by
$\varepsilon v_i e_i$ changes the corresponding arm length by
$O(\varepsilon t^{-2})$. Transport bounds the change in $D$, and
the same bound follows from the formula for $D_*$. These bounds
control the negative part uniformly after division by $\varepsilon t$.
On a face the same estimates apply in the tangential variables.
Fatou's lemma now proves Equation~\eqref{eq:src11}. Each limiting
integral has finite negative part since $g_i$ is semiconvex near
zero and takes its minimum there.

The right derivatives $g_i'(0+)$ exist and are nonnegative. A
positive derivative would, by semiconvexity, give a positive lower
bound for $g_i'$ near zero and hence an infinite integral in
Equation~\eqref{eq:src11}. Thus every $g_i'(0+)$ vanishes. The
directional derivative of $\Delta$ at $v$ is sublinear by
semiconvexity and nonnegative by minimality. Since it vanishes on
the positive coordinate vectors, it vanishes on their positive span.
In particular $K=0$. Integrating by parts gives
\[
       -J\ \ge\ \sum_{i=1}^k v_i
                    \int_0^\infty\frac{g_i(t)}{t^2}\,dt\ \ge\ 0.
\]
The boundary terms vanish because $g_i(t)=o(t)$ at zero and
Equation~\eqref{eq:defectbound} holds at infinity. As $J\ge0$,
continuity gives $g_i(t)=0$ for every $i$ and $t\ge0$. Applying
the same argument successively along the coordinate rays reaches
every point above $v$ coordinatewise. Any positive triple, or any
point of the same open face, can be dilated into this region.
Homogeneity gives zero defect throughout the interior or the face
in which we started.

It remains to move a face zero into the interior. Suppose $v_1=0$
and $\Delta(v)=0$. At $v^{(r)}=v+re_1$, tangential dominated
convergence gives $J_2(v^{(r)})\to J_2(v)=0$. By the uniform
Equation~\eqref{eq:src10} and continuity of $\Delta$,
$J_1(v^{(r)})\to0$ as well. Although we have only tangential
Lipschitz bounds at the face, transport shows that
$\Delta-h=D-\sum_i a_i$ is nondecreasing in $v_1$. Since the normal
derivative of $h$ stays bounded near this point,
\[
 \Delta(v^{(r)}+te_1)-\Delta(v^{(r)})\ge-Ct \qquad(0<t\le1).
\]
Fatou's lemma in the normal integral now yields
$\int_0^\infty\Delta(v+te_1)\,dt/t\le0$. The continuous nonnegative
integrand vanishes at every $t>0$, giving the interior zero required
above. In particular a homogeneous tripod with contact in either
regime has $D(1,1,1)=3-1/3=8/3$, the equal-arm value used in
Section~\ref{sec:boundary}.
\end{proof}

\subsection{The three wire lengths at the maximal price}

The unconditional inequality of Proposition~\ref{prop:tripod} gives a
geometric proof of the known square upper bound and identifies the wire
prices when equality holds. Work at scale one on three sites, so that a
label shared by all three has rate one. Write the three wire lengths as
\(a_i+a_j\) for \(i\ne j\). The triangle inequalities give \(a_i\ge0\).
Over an algebraically closed field the Proposition and the
arithmetic--harmonic mean inequality give
\begin{equation}\label{eq:arm-perimeter}
  1\ \ge\ \sum_i a_i-\Big(\sum_i a_i^{-1}\Big)^{-1}
     \ \ge\ \frac89\sum_i a_i .
\end{equation}
The reciprocal term is zero if an arm vanishes.
The joint values of the three wires are preserved on passing to an
algebraic closure (Section~\ref{sec:prices-background}), so the same
bound holds over every field. Since \(T_n\) consists of one
\(n\)-dimensional wire on each side, every spectral point assigns it
value \(n^{2\sum_i a_i}\le n^{9/4}\). Strassen's maximum formula for
asymptotic rank~\cite{spectrum} gives the upper bound. Fixed-size
recursion as in Lemma~\ref{lem:specialize} then gives \(\omega_F\le9/4\).

\begin{corollary}\label{cor:extremal-price}
Fix \(n\ge2\) and a field \(F\). At every scale-one spectral point
maximizing the value of \(T_n\), all three wire lengths are \(3/4\).
There exists such a point, and it assigns \(T_{a,b,c}\) the value
\((abc)^{3/4}\) for every choice of positive integers \(a,b,c\).
\end{corollary}

\begin{proof}
The spectral maximum is attained. Corollary~\ref{cor:tight} identifies
its value as \(n^{9/4}\), so \(\sum_i a_i=9/8\). In
\eqref{eq:arm-perimeter} both inequalities must be equalities. All three
arms are positive, since a zero arm would give \(1\ge9/8\).
Equality of the arithmetic and harmonic means gives \(a_i=3/8\) for
each \(i\). Each wire therefore has rate \(3/4\). The rectangular tensor
has one wire of each dimension \(a,b,c\), and their values multiply.
\end{proof}
